% ch3_d 第3章 §3.4–§3.5 (书页 101–111 / p116–p126)
%--C-TAIL--
其中 $K(\pmb{x}) = n^{2}\langle \delta\theta(\pmb{x})\delta\theta(0)\rangle$。我们注意到，最后一项可以改写为

% ---- 书页 101 (p116) ----
在超流相中，只有相位涨落在长距离下才是重要的。令 $\varphi_c(\pmb{x}) = \varphi_0\mathrm{e}^{\mathrm{i}\theta(\pmb{x})}$，我们便得到一个 XY 模型
\begin{equation*}
S_{\mathrm{eff}} = \int \mathrm{d}^{d}\pmb{x}\ \frac{\eta}{2}(\partial_{\pmb{x}}\theta)^{2}
\tag{3.4.2}
\end{equation*}
其中
\begin{equation*}
\eta = \frac{\varphi_{0}^{2}}{mT}
\tag{3.4.3}
\end{equation*}
我们注意到，能量由 $TS_{\mathrm{eff}}$ 给出。因此，$T\eta$ 的数值决定了相位扭转 $\partial_{\pmb{x}}\theta \neq 0$ 的能量代价。这里 $T\eta$ 称为 XY 模型的相位劲度（phase rigidity）。

当然，以上关于对称性破缺的讨论都基于经典图像，其中我们忽略了 $\varphi$ 和 $\theta$ 的涨落。（现在这些涨落对应于热涨落。）我们同样可以问：对称性破缺态能否在热涨落下存活？与上一节一样，我们可以用 XY 模型来处理这个问题。利用那些结果，我们发现，对于 $d>2$，热涨落并不总是破坏长程序。

对于 $d<2$，热涨落总是破坏长程序，把它变为短程序。这是因为，对于 $d<2$，$\theta$ 的关联随 $|\pmb{x}|$ 呈幂次发散：
\begin{equation*}
\begin{aligned}
&\ \langle \theta(\pmb{x})\theta(0)\rangle - \langle \theta(l)\theta(0)\rangle
 = C_{d}\eta^{-1}(L^{2-d} - |\pmb{x}|^{2-d}) - C_{d}\eta^{-1}(L^{2-d} - l^{2-d}) \\
&\ = -C_{d}\eta^{-1}|\pmb{x}|^{2-d}
\end{aligned}
\end{equation*}
这里用到了傅里叶变换
\begin{equation*}
\int \mathrm{d}^{d}\pmb{k}\ \frac{\mathrm{e}^{\mathrm{i}\pmb{k}\cdot\pmb{x}}}{\eta|\pmb{k}|^{2}} = C_{d}\eta^{-1}(L^{2-d} - |\pmb{x}|^{2-d})
\end{equation*}
因此，
\begin{equation*}
\left\langle \mathrm{e}^{\mathrm{i}n\theta(\pmb{x})}\mathrm{e}^{-\mathrm{i}n\theta(0)} \right\rangle
= \mathrm{e}^{-\frac{n^{2}}{2}\langle \theta(\pmb{x})-\theta(0)\rangle^{2}}
= \mathrm{e}^{n^{2}\langle \theta(\pmb{x})\theta(0)\rangle - n^{2}\langle \theta(l)\theta(0)\rangle}
= \mathrm{e}^{-C_{d}n^{2}\eta^{-1}|\pmb{x}|^{2-d}}
\end{equation*}
当 $n=1$ 时，衰减长度为 $(\eta/C_{d})^{1/(2-d)}$。

对于 $d=2$，$\theta$ 关联具有对数发散 $\langle \theta(\pmb{x})\theta(0)\rangle - \langle \theta(l)\theta(0)\rangle = -\frac{1}{2\pi\eta}\ln(|\pmb{x}|/l)$，我们有以下的代数长程序（algebraic long-range order）：
\begin{equation*}
\left\langle \mathrm{e}^{\mathrm{i}n\theta(\pmb{x})}\mathrm{e}^{-\mathrm{i}n\theta(0)} \right\rangle \sim |\pmb{x}|^{-n^{2}/2\pi\eta}
\tag{3.4.4}
\end{equation*}

% ---- 书页 102 (p117) ----
\subsection{Kosterlitz--Thouless 相变}
\begin{keypoints}
\item 我们已经看到，依赖空间的 $\theta$ 涨落会把长程序变为代数长程序。这里我们将看到，一类新的涨落——涡旋涨落——能够把代数长程序变为短程序。
\end{keypoints}

上面的 $d=2$ 结果并不完全正确。当 $\eta$ 低于临界值 $\eta_{c}$（或温度高于临界温度 $T_{KT}$）时，代数长程序无法存活，只剩下短程关联（Kosterlitz and Thouless, 1973）。为了理解这一现象，我们需要把涡旋涨落包括进来。一个涡旋位形由下式给出
\begin{equation*}
\varphi_{c}(x,y) = f(r)\mathrm{e}^{\mathrm{i}\phi}, \qquad f(0)=0, \qquad f(\infty)=\varphi_{0}
\end{equation*}
其中 $x = r\cos(\phi)$，$y = r\sin(\phi)$。可以证明 $|\partial_{\pmb{x}}\theta| = 1/r$。因此，单个涡旋的作用量为
\begin{equation*}
S_{v} = \int \mathrm{d}^{2}r\ \frac{1}{2}\eta\frac{1}{r^{2}} + S_{c}
= \int \pi\,\mathrm{d}(r^{2})\ \frac{1}{2}\eta\frac{1}{r^{2}} + S_{c}
= h\ln\frac{L}{l} + S_{c}
\end{equation*}
其中
\begin{equation*}
h = \eta\pi,
\end{equation*}
$l$ 是短距离截断尺度，$S_{c}$ 是芯部作用量（core action，即来自涡旋芯内部的贡献，在那里 $f(r)$ 开始从 $\varphi_{0}$ 变为 0）。相距 $r$ 的一个涡旋与一个反涡旋（anti-vortex）之间的相互作用为 $2h\ln\frac{r}{l}$。

这里的问题与 2.4.1 节和 2.4.3 节讨论的瞬子气体问题非常相似。为了计算配分函数，我们必须包括涡旋的贡献。在固定涡旋位形 $\varphi_{c} = \mathrm{e}^{\mathrm{i}\theta_{c}}$ 之下，使用 XY 模型，配分函数具有如下形式
\begin{equation*}
Z = \int \mathcal{D}\delta\theta\ \mathrm{e}^{-\int \mathrm{d}^{2}\pmb{x}\ \frac{\eta}{2}\left(\partial_{\pmb{x}}(\delta\theta+\theta_{c})\right)^{2}}
\end{equation*}
（译注：原文导数下标印作 $n$，按上下文应为 $\pmb{x}$。）

由于 XY 模型是二次的，所得的配分函数取简单的因子化形式 $Z = Z_{0}\mathrm{e}^{-S_{\mathrm{eff}}(\theta_{c})}$，其中 $Z_{0}$ 是没有涡旋时的配分函数。由于 $S_{\mathrm{eff}}(\theta_{c})$ 由涡旋之间的对数相互作用给出，总配分函数具有如下形式
\begin{equation*}
Z = Z_{0} \sum_{n} \frac{1}{n!n!} \int \prod_{i=1}^{2n} \mathrm{d}^{2}\pmb{r}_{i}\ \mathrm{e}^{-2nS_{c}}\ \mathrm{e}^{\sum_{i<j}^{2n} 2hq_{i}q_{j}\ln\frac{r_{ij}}{l}}
\tag{3.4.5}
\end{equation*}
求和中的每一项包含位于 $\pmb{r}_{1},\ldots,\pmb{r}_{n}$ 的 $n$ 个涡旋和位于 $\pmb{r}_{n+1},\ldots,\pmb{r}_{2n}$ 的 $n$ 个反涡旋。此外，$q_{i} = \pm$ 取决于第 $i$ 个涡旋是涡旋还是反涡旋，而 $r_{ij}$ 是第 $i$ 个与第 $j$ 个（反）涡旋之间的距离。

% ---- 书页 103 (p118) ----
%FIG{3.8}: {$S_{\mathrm{eff}}$ 作为 $l/l_{n}$ 的函数。四个面板分别对应 $S_{c}\gg 1,\ h<2$；$S_{c}\ll -1,\ h<2$；$S_{c}\gg 1,\ h>2$；$S_{c}\ll -1,\ h>2$。}
为了理解涡旋的效应，让我们估计 $n$ 个涡旋与 $n$ 个反涡旋的配分函数，即 $Z = \mathrm{e}^{-S_{\mathrm{eff}}}$，其中
\begin{equation*}
S_{\mathrm{eff}} \sim 2n\ln n + n\left(2h\ln\frac{l_{n}}{l} + 2S_{c}\right) - 2n\ln\frac{L^{2}}{l^{2}}
= L^{2}\frac{2}{l_{n}^{2}}(h-2)\ln\frac{l_{n}\mathrm{e}^{S_{c}/(h-2)}}{l}
\tag{3.4.6}
\end{equation*}
其中 $L$ 是系统的线性尺寸，$l_{n} = L/\sqrt{n}$ 是涡旋之间的平均间距。这里 $l_{n}$ 总是大于截断，即 $l_{n} > l$。式 (3.4.6) 的第一项来自 $1/(n!)^{2}$。第二项是 $n$ 对涡旋--反涡旋的作用量，每对占据大小 $\sim l_{n}$ 的一块面积。它包含两项贡献：涡旋相互作用 $2h\ln\frac{l_{n}}{l}$ 和芯部作用量 $2S_{c}$。第三项来自对涡旋位置的积分 $\int \prod_{i=1}^{2n}\mathrm{d}^{2}\pmb{r}_{i}$（它就是熵）。涡旋气体的行为由涡旋作用量（倾向于更少的涡旋）与熵（倾向于更多的涡旋）之间的竞争控制。

涡旋的数目 $n$ 可以这样简单地得到：在 $l_{n} > l$ 的范围内对 $l_{n}$ 极小化 $S_{\mathrm{eff}}$。从图 3.8 中 $S_{\mathrm{eff}}$ 的行为我们看到，当 $S_{c} \ll -1$（即产生涡旋的代价很低）时，$S_{\mathrm{eff}}$ 在边界 $l_{n} \sim l$ 处取极小。在这种情况下，涡旋涨落是重要的，它把式 (3.4.4) 中的代数长程序变为短程序：
\begin{equation*}
\left\langle \mathrm{e}^{\mathrm{i}\theta(\pmb{x})}\mathrm{e}^{-\mathrm{i}\theta(0)} \right\rangle \sim \mathrm{e}^{-|\pmb{x}|/\xi}
\end{equation*}
关联长度为 $\xi \sim l_{n} \sim l$。

当 $S_{c} \gg 1$ 且 $h = \pi\eta < 2$ 时，$S_{\mathrm{eff}}$ 在 $l_{n} \sim \mathrm{e}^{S_{c}/(2-h)}l$ 处取极小。此时涡旋密度仍然是有限的，这同样把代数长程序变为

% ---- 书页 104 (p119) ----
%FIG{3.9}: {二维 XY 模型的相图。横轴为 $\eta$，在 $\eta_{c}$ 之左为短程序，之右为代数长程序。}
短程序。关联长度为
\begin{equation*}
\xi \sim l_{n} \sim l\,\mathrm{e}^{S_{c}/(2-h)}
\tag{3.4.7}
\end{equation*}
当 $S_{c} \gg 1$ 且 $h = \pi\eta > 2$ 时，$S_{\mathrm{eff}}$ 在 $l_{n} = \infty$ 处取极小，涡旋密度为零。在这种情况下，涡旋涨落在长距离下并不重要，代数长程序可以在涡旋涨落下存活。

从上面的讨论我们看到，在 $S_{c} \gg 1$ 极限下，随着我们改变 $\eta$，二维玻色系统经历一个有限温度相变。该相变把代数长程序变为短程序。相变处的临界 $\eta$ 为 $\eta_{c} = 2/\pi$。图 3.9 中的相图总结了我们在 $S_{c} \gg 1$ 极限下的结果。$\eta_{c}$ 处的相变就是著名的 Kosterlitz--Thouless（KT）相变。利用 $\eta$ 与 $T$ 之间的关系 (3.4.3)，我们求得临界温度
\begin{equation*}
T_{KT} = \frac{\pi\varphi_{0}}{2m}
\end{equation*}
我们想指出，KT 相变不改变任何对称性，因而是 Landau 关于相与相变的对称性破缺理论的一个反例。

利用 $h = \pi\eta = \pi\varphi_{0}^{2}/mT$，式 (3.4.7) 可以改写为
\begin{equation*}
\xi(T) \sim l\,\mathrm{e}^{E_{\xi}/(T-T_{KT})}
\tag{3.4.8}
\end{equation*}

\begin{smallnote}
在 $T_{KT}$ 附近，其中 $E_{\xi}$ 是某个能量尺度。这里我们想指出，当 $T \to T_{KT}$ 时式 (3.4.8) 并不正确。这个不正确的结果来自一个不正确的假设，即相位劲度 $\eta$ 不受涡旋存在的影响。事实上，涡旋气体会修改 $\eta$ 和 $h$ 的有效值。这是因为涡旋--反涡旋对可以被相位扭转 $\partial_{\pmb{x}}\theta$ 极化，从而释放部分应变并降低有效相位劲度。采用有效值 $h^{*}$，式 (3.4.7) 应当改写为
\begin{equation*}
\xi \sim l_{n} \sim l\,\mathrm{e}^{S_{c}/(2-h^{*})}
\tag{3.4.9}
\end{equation*}
$h^{*}(T)$ 的温度依赖性比 $h(T)$ 的更为复杂。不过，按定义 $h^{*}(T_{KT}) = 2$。基于一个重整化群计算（见 3.5.6 节），我们发现
\begin{equation*}
2 - h^{*}(T_{KT}) \propto (T - T_{KT})^{1/2}
\end{equation*}
（译注：原文括号内印作 $T_{TK}$，应为 $T_{KT}$。）
由此得到
\begin{equation*}
\xi(T) \sim l\,\mathrm{e}^{\left(E_{\xi}/(T-T_{KT})\right)^{1/2}}
\tag{3.4.10}
\end{equation*}
\end{smallnote}

% ---- 书页 105 (p120) ----
\subsection*{问题 3.4.1}
注意到，在 $T_{c}$ 附近，式 (3.4.1) 中的作用量 $S_{\mathrm{eff}}$ 可以近似为
\begin{equation*}
S_{\mathrm{eff}} = \beta\int \mathrm{d}^{d}\pmb{x}\ \left[ \frac{1}{2m^{*}}|\partial_{\pmb{x}}\varphi_{c}|^{2} + a(T-T_{c})|\varphi_{c}|^{2} + b|\varphi_{c}|^{4} \right] + O(|\varphi_{c}|^{6})
\end{equation*}
因为在 $T = T_{c}$ 处的临界点（或相变点）附近序参量 $\varphi_{c}$ 很小。在用平均场（或半经典）方法处理相变和临界点时，我们首先求出使作用量取极小的平均场解，然后假设平均场解附近的涨落很小，把作用量展开到涨落的二次阶。$S_{\mathrm{eff}}$ 的二次近似可以用来计算各种关联函数。

\begin{enumerate}
\item 用平均场方法计算临界点处 $\langle \varphi_{c}(\pmb{x})\varphi_{c}^{*}(0)\rangle \sim 1/|\pmb{x}|^{\gamma}$ 中的衰减指数 $\gamma$。
\item 上述结果并非总是成立，因为经典理论可能会失效。重复 3.3.8 节末尾的讨论（即把 $S_{\mathrm{eff}}$ 写成 $g^{-1}\tilde{S}$ 的形式，$\tilde{S}$ 无量纲），看看平均场方法何时能正确描述临界点、何时临界点由强涨落控制；也就是说，求出上临界点 $d_{c}$。
\item 这里我们想引入相关微扰与不相关微扰的概念。我们知道，在上临界维数以上，经典理论能正确描述相变处的临界点。现在我们在有效作用量 $S_{\mathrm{eff}}$ 中加一个微扰 $\beta\int \mathrm{d}^{d}x\ c|\varphi|^{\sigma}$。如果该微扰是重要的并且修改临界指数，我们就称它为相关微扰（relevant perturbation）；如果该微扰在临界点附近变得任意小，我们就称它为不相关微扰（irrelevant perturbation）。利用你在上面找到的同样的标度关系，考察微扰 $\beta\int \mathrm{d}^{d}x\ c|\varphi|^{\sigma}$ 如何修改标度后的作用量 $\tilde{S}$。确定 $\sigma$ 在什么范围内该微扰是相关的，在什么范围内是不相关的。
\end{enumerate}

\section{重整化群}

\subsection{相关微扰与不相关微扰}
\begin{keypoints}
\item 相关微扰改变系统的长距离（或低能）行为，而不相关微扰不改变。
\item 我们可以利用微扰的标度量纲来判断该微扰是相关的还是不相关的。
\end{keypoints}

在上面关于 KT 相变的讨论中我们注意到，当 $\mathrm{e}^{-S_{c}} \ll 1$ 时，涡旋涨落只是“小微扰”。然而，如果 $h < 2$，那么无论 $\mathrm{e}^{-S_{c}}$ 多么小，涡旋涨落总会破坏 $\langle \mathrm{e}^{\mathrm{i}\theta(\pmb{x})}\mathrm{e}^{-\mathrm{i}\theta(0)}\rangle$ 的代数长程关联。因此，当 $h<2$ 时，把涡旋涨落包括进来这一微扰被称为相关微扰。当 $h>2$ 时，该微扰被称为不相关微扰；而当 $h=2$ 时，该微扰被称为边缘微扰（marginal perturbation）。下面我们想在更一般的框架中讨论相关/不相关/边缘微扰。

% ---- 书页 106 (p121) ----
考虑一个由作用量
\begin{equation*}
S = S_{0} + \int \mathrm{d}^{d}x\ aO(x)
\end{equation*}
描写的理论，其中 $aO$ 是一个微扰。我们假设 $S_{0}$ 具有 $Z_{2}$ 对称性，且在 $Z_{2}$ 变换下 $O \to -O$。于是，当 $a=0$ 时 $\langle O\rangle = 0$。我们还假设，对于大的 $x$，在 $a=0$ 时
\begin{equation*}
\langle O(x)O(0)\rangle = \frac{1}{|x|^{2h}}
\tag{3.5.1}
\end{equation*}
这里 $h$ 称为算符 $O$ 的标度量纲（scaling dimension，$1/x$ 的标度量纲定义为 1）。式 (3.5.1) 同时也定义了算符 $O$ 的归一化。

在二阶微扰下，配分函数由下式给出
\begin{equation*}
Z = Z_{0} \int \mathrm{d}^{d}x\,\mathrm{d}^{d}y\ a^{2}\langle O(x)O(y)\rangle
\end{equation*}
其中 $Z_{0}$ 是零阶配分函数。我们看到，二阶微扰对有效作用量的改变为
\begin{equation*}
\Delta S_{\mathrm{eff}} = -\ln Z + \ln Z_{0} = -2\ln g + 2h\ln L - 2d\ln L
\end{equation*}
（译注：原文 $\ln g$ 按上下文应为 $\ln a$。）

注意到，当 $h < d$ 且 $L > \xi = a^{-1/(d-h)}$ 时，我们有 $\Delta S_{\mathrm{eff}} < 0$。因此，系统倾向于有两个 $O(x)$ 插入。当 $L \gg \xi$ 时，系统希望在每个 $\xi^{d}$ 体积内有两个 $O(x)$ 插入。我们看到，如果我们感兴趣的是 $\xi$ 之外长度尺度上的关联函数，那么该微扰总是重要的。我们得出结论：如果 $O(x)$ 的标度量纲小于 $d$，则微扰 $\int \mathrm{d}^{d}x\ O(x)$ 是相关的。此时 $O(x)$ 称为相关算符。如果 $O(x)$ 的标度量纲大于（或等于）$d$，则 $O(x)$ 称为不相关（边缘）算符。记住这个结果的一个简单办法是：注意若 $\int \mathrm{d}^{d}x\ O(x)$ 的量纲小于零，则微扰 $\int \mathrm{d}^{d}x\ O(x)$ 是相关的。

标度量纲的概念还使我们能够用量纲分析来估计有限微扰 $aO$ 所诱导的 $\langle O\rangle$。由于 $\delta S = \int \mathrm{d}^{d}x\ aO$ 的标度量纲按定义为零，系数 $a$ 具有标度量纲 $[a] = d - [O] = d - h$。当 $aO$ 是不相关微扰（即 $h > d$）时，诱导的 $\langle O\rangle$ 正比于 $a$。我们有
\begin{equation*}
\langle O\rangle \sim a\,l^{d-2h}
\end{equation*}
其中 $l$ 是短距离截断。当 $aO$ 是相关微扰（即 $h < d$）时，诱导的 $\langle O\rangle$ 大于 $a\,l^{d-2h}$。通过匹配标度量纲，

% ---- 书页 107 (p122) ----
我们求得
\begin{equation*}
\langle O\rangle \sim a^{h/(d-h)} = a\xi^{d-2h}.
\tag{3.5.2}
\end{equation*}

\subsection*{问题 3.5.1}
有效作用量
\begin{equation*}
S_{\mathrm{eff}} = \beta\int \mathrm{d}^{d}\pmb{x}\ \frac{1}{2m^{*}}|\partial_{\pmb{x}}\varphi|^{2}
\end{equation*}
描写一个临界点。计算 $\varphi$、$\varphi^{2}$、$|\varphi|^{2}$ 和 $|\varphi|^{4}$ 的标度量纲。证明：在某个空间维度 $d_{0}$ 以下，微扰 $\int \mathrm{d}^{d}\pmb{x}\ b|\varphi|^{4}$ 成为相关微扰。求出 $d_{0}$ 的值，并解释为什么 $d_{0}$ 等于
\begin{equation*}
S_{\mathrm{eff}} = \beta\int \mathrm{d}^{d}\pmb{x}\ \left[ \frac{1}{2m^{*}}|\partial_{\pmb{x}}\varphi|^{2} + a(T-T_{c})|\varphi|^{2} + b|\varphi|^{4} \right]
\end{equation*}
的上临界维数 $d_{c}$。

\subsection{二维 XY 模型与二维时钟模型之间的对偶}
\begin{keypoints}
\item 二维 XY 模型中的涡旋可以看作粒子。描写这些粒子的场论是二维时钟模型（clock model）。
\end{keypoints}

为了更仔细地研究 XY 模型的涡旋涨落，我们想把二维 XY 模型映射到 $Z_{1}$ 二维时钟模型。一个一般的 $Z_{n}$ 时钟模型定义为
\begin{equation*}
S = \int \mathrm{d}^{2}\pmb{x}\ \left( \frac{\kappa}{2}(\partial_{\pmb{x}}\theta)^{2} - g\cos(n\theta) \right)
\tag{3.5.3}
\end{equation*}
当 $g = 0$ 时，时钟模型就是有限温度下的 XY 模型。作用量是能量除以温度：$S = \beta E$。$g\cos(n\theta)$ 项（显式地）破坏 $U(1)$ 转动对称性。如果我们把 $(\cos(\theta), \sin(\theta)) = (S_{x}, S_{y})$ 看作自旋的两个分量，那么对于 $n = 1$，$g\cos(\theta)$ 项就是由 $S_{x}$ 方向的磁场诱导的项。对于一般的 $n$，时钟模型具有 $Z_{n}$ 对称性：$\theta \to \theta + \frac{2\pi}{n}$。

为了证明对偶关系，我们考虑式 (3.5.3) 在 $n = 1$ 时的配分函数：
\begin{equation*}
\begin{aligned}
Z &= \int \mathcal{D}\theta\ \sum_{k} \frac{1}{k!}\left( g\int \mathrm{d}^{2}x\ \frac{\mathrm{e}^{\mathrm{i}\theta} + \mathrm{e}^{-\mathrm{i}\theta}}{2} \right)^{k} \mathrm{e}^{-\int \mathrm{d}^{2}\pmb{x}\ \frac{\kappa}{2}(\partial_{\pmb{x}}\theta)^{2}} \\
&= Z_{0} \sum_{k} \frac{1}{k!k!} \int \prod_{i=1}^{2k} \mathrm{d}^{2}\pmb{r}_{i}\ \mathrm{e}^{-2kS_{c}}\ \mathrm{e}^{\sum_{i<j}^{2k} 2hq_{i}q_{j}\ln\frac{r_{ij}}{l}}
\end{aligned}
\tag{3.5.4}
\end{equation*}
其中 $Z_{0}$ 是 $S = \int \mathrm{d}^{2}\pmb{x}\ \frac{\kappa}{2}(\partial_{\pmb{x}}\theta)^{2}$ 的配分函数。求和中的每一项来自关联 $\langle \mathrm{e}^{\mathrm{i}\theta(\pmb{r}_{1})}\allowbreak\ldots\allowbreak\mathrm{e}^{\mathrm{i}\theta(\pmb{r}_{k})}\allowbreak\mathrm{e}^{-\mathrm{i}\theta(\pmb{r}_{k+1})}\allowbreak\ldots\allowbreak\mathrm{e}^{-\mathrm{i}\theta(\pmb{r}_{2k})}\rangle$。此外，

% ---- 书页 108 (p123) ----
$q_{i} = 1$（当 $1 \leqslant i \leqslant k$），$q_{i} = -1$（当 $k+1 \leqslant i \leqslant 2k$），$\mathrm{e}^{-S_{c}} = g/2$，以及
\begin{equation*}
h = \frac{1}{4\pi\kappa}
\end{equation*}
（译注：原文 $q_{i} = -1$ 的条件印作 $k+1 \leqslant i \leqslant 2n$，应为 $2k$。）

如果 $\frac{1}{4\pi\kappa} = \pi\eta$，则式 (3.5.4) 与 XY 模型 (3.4.2) 的配分函数 (3.4.5) 完全相同。所以，当 $\frac{1}{2\pi\kappa} = 2\pi\eta$ 时，$Z_{1}$ 时钟模型 (3.5.3) 等价于（带涡旋的）XY 模型 (3.4.2)。XY 模型中的涡旋被映射到时钟模型中的 $\mathrm{e}^{\mathrm{i}\theta(\pmb{x})}$。类似地，时钟模型中的涡旋被映射到 XY 模型中的 $\mathrm{e}^{\mathrm{i}\theta(\pmb{x})}$。时钟模型中的涡旋具有标度量纲 $\pi\kappa$，而 XY 模型中的 $\mathrm{e}^{\mathrm{i}\theta(\pmb{x})}$ 算符具有标度量纲 $1/4\pi\eta$。关系 $\frac{1}{2\pi\kappa} = 2\pi\eta$ 保证了这两个标度量纲彼此一致。

我们知道，XY 模型中的 $U(1)$ 对称性不允许 $\mathrm{e}^{\mathrm{i}\theta}$ 项出现在作用量中。利用上述映射我们看到，相应的时钟模型必须不允许涡旋涨落。在时钟模型中允许涡旋涨落，对应于在对偶的 XY 模型中显式破坏 $U(1)$ 对称性。我们看到，时钟模型有两种不同的类型：允许涡旋涨落的一种和不允许涡旋涨落的一种。由于相应的对偶模型具有不同的对称性，这两种时钟模型具有非常不同的性质。我们把允许涡旋涨落的时钟模型称为紧致时钟模型（compact clock model），把不允许涡旋涨落的称为非紧致时钟模型（non-compact clock model）。带涡旋的 XY 模型被映射到一个 $Z_{1}$ 非紧致时钟模型。这样的映射使我们能够通过研究相应非紧致时钟模型中的相变来研究 XY 模型中的 KT 相变。

\subsection{时钟模型的物理性质}
\begin{keypoints}
\item 除非指定短距离截断，一个场论模型是没有良定义的。
\item 包含强涡旋涨落的 Ginzburg--Landau 理论不能描述非紧致时钟模型中的相变。
\end{keypoints}

本节中，我们将讨论一般的 $Z_{n}$ 时钟模型中可能的相变。当 $g$ 大时，场 $\theta$ 被势 $-g\cos(n\theta)$ 的某个极小值俘获。我们相信，此时模型处于一个自发破缺 $Z_{n}$ 对称性的相。当 $\kappa$ 和 $g$ 都小时，$\theta$ 的涨落很强。我们预期模型将处于 $Z_{n}$ 对称相。

尽管听起来相当合理，上述说法其实并没有真正的意义。这是因为 $g$ 具有量纲，谈论 $g$ 有多大是没有意义的。更糟的是，$g$ 是模型中唯一具有非平凡量纲的参量，所以我们无法构造一个无量纲组合来判断 $g$ 有多大。

% ---- 书页 109 (p124) ----
为了以物理的方式理解 $g\cos(n\theta)$ 项的重要性，我们想问 $\mathrm{e}^{\mathrm{i}n\theta}$ 算符有多大。回答这个问题的一个物理办法，是考察 XY 模型 $S = \int \mathrm{d}^{2}\pmb{x}\ \frac{\kappa}{2}(\partial_{\pmb{x}}\theta)^{2}$ 中 $\mathrm{e}^{\mathrm{i}n\theta}$ 的关联。该关联由下式给出（见式 (3.3.25)）
\begin{equation*}
\left\langle \mathrm{e}^{\mathrm{i}n\theta(\pmb{x})}\mathrm{e}^{-\mathrm{i}n\theta(0)} \right\rangle = (l/|\pmb{x}|)^{n^{2}/2\pi\kappa}
\tag{3.5.5}
\end{equation*}
一个大的意外是，该关联依赖于短距离截断 $l$。这样一来，算符 $\mathrm{e}^{\mathrm{i}n\theta}$ 的大小（或重要性）在未指定截断 $l$ 之前甚至没有良定义。这说明了这样一点：\emph{要有一个良定义的场论，我们必须指定一个短距离截断 $l$}。为了强调这一点，我们想把 $l$ 依赖性显式地写出来，把作用量写成
\begin{equation*}
S = \int \mathrm{d}^{2}\pmb{x}\ \left( \frac{\kappa_{l}}{2}(\partial_{\pmb{x}}\theta_{l})^{2} - g_{l}\cos(n\theta_{l}) \right)
\tag{3.5.6}
\end{equation*}
短距离截断是这样引入的：要求 $\theta_{l}$ 场不包含任何波长短于 $l$ 的涨落：
\begin{equation*}
\theta_{l}(\pmb{x}) = \int_{|\pmb{k}|<2\pi/l} \mathrm{d}^{2}\pmb{k}\ \theta_{\pmb{k}}\mathrm{e}^{\mathrm{i}\pmb{x}\cdot\pmb{k}}
\end{equation*}
我们看到，良定义的时钟模型 (3.5.6) 包含三个参量 $\kappa_{l}$、$g_{l}$ 和 $l$。所以，时钟模型实际上包含两个无量纲参量 $\kappa_{l}$ 和 $\bar{g}_{l} = g_{l}l^{2}$。

现在我们可以做有意义的陈述了。当 $\bar{g}_{l} \gg 1$ 时，我们相信模型处于自发破缺 $Z_{n}$ 对称性的相。当 $\kappa_{l}$ 和 $\bar{g}_{l}$ 都远小于 1 时，我们预期模型处于 $Z_{n}$ 对称相。

一个非平凡的极限是 $\kappa_{l} \gg 1$ 且 $\bar{g}_{l} \ll 1$。此时模型是处于 $Z_{n}$ 对称相还是 $Z_{n}$ 对称性破缺相？相关/不相关微扰的概念对回答这个问题非常有帮助。如果我们把 $g_{l}\cos(n\theta)$ 项当作 XY 模型的微扰，那么由式 (3.5.5)，XY 模型中 $\mathrm{e}^{\mathrm{i}n\theta}$ 的标度量纲为 $[\mathrm{e}^{\mathrm{i}n\theta}] = n^{2}/4\pi\kappa_{l}$。因此，$g_{l}\cos(n\theta)$ 项在 $n^{2}/4\pi\kappa_{l} < 2$ 时是相关的，在 $n^{2}/4\pi\kappa_{l} > 2$ 时是不相关的。

这个结果是合理的。当 $\kappa_{l}$ 小时，$\theta$ 的涨落强。这使得 $g_{l}\cos(n\theta)$ 项平均为零，效果减弱。因此微扰 $g_{l}\cos(n\theta)$ 是不相关的。当 $g_{l}\cos(n\theta)$ 不相关且 $\bar{g}_{l}$ 小时，我们在计算长程关联时可以丢掉 $g_{l}\cos(n\theta)$ 项。这表明，在 $\kappa_{l}$ 和 $\bar{g}_{l}$ 都小时，长距离下我们不仅有 $Z_{n}$ 对称性，还有完整的 $U(1)$ 对称性。

当 $\kappa$ 大时，$\theta$ 的涨落弱。这使得 $g_{l}\cos(n\theta)$ 项成为相关微扰。此时无论 $\bar{g}_{l}$ 多么小，$g_{l}\cos(n\theta)$ 项的效应都会在长距离下变得重要。

% ---- 书页 110 (p125) ----
这样，我们预期系统被俘获在势项的 $n$ 个极小之一中，$Z_{n}$ 对称性自发破缺，即使 $\bar{g}_{l}$ 很小。

在认识到时钟模型可以有 $Z_{n}$ 对称性破缺相和 $Z_{n}$ 对称相之后，下一个自然的问题是：这两相之间如何相互转化？理解这种相变的一个办法，是引入复序参量 $\varphi(\pmb{x}) = \langle \mathrm{e}^{\mathrm{i}\theta(\pmb{x})}\rangle$，并写下相变的 Ginzburg--Landau 有效理论
\begin{equation*}
S_{GL} = \int \mathrm{d}^{2}\pmb{x}\ \left[ \frac{\gamma}{2}(\partial_{\pmb{x}}\varphi)^{2} + a|\varphi|^{2} + b|\varphi|^{4} - c\,\mathrm{Re}\,\varphi^{n} \right]
\tag{3.5.7}
\end{equation*}
注意到，$c\,\mathrm{Re}\,\varphi^{n}$ 项（显式地）把 $U(1)$ 对称性破缺到 $Z_{n}$。当 $n > 1$ 时，随着 $a$ 从正值变为负值，Ginzburg--Landau 理论描述了一个对称性破缺相变。

当 $n = 1$ 时，Ginzburg--Landau 理论不包含相变，因为没有对称性破缺。这似乎暗示 $Z_{1}$ 时钟模型不包含相变，而相应的 XY 模型不包含 KT 相变。

那么问题出在哪里？在 Ginzburg--Landau 理论中，序参量在相变点附近有很强的振幅涨落。$\varphi$ 的一个典型位形包含许多 $\varphi = 0$ 的点，所以存在很强的涡旋涨落。Ginzburg--Landau 理论描述的是紧致时钟模型中的相变，它不适用于非紧致时钟模型。

\subsection{非紧致时钟模型的重整化群方法}
\begin{keypoints}
\item 通过跑动耦合常数（running coupling constants）的概念，重整化群（RG）方法使我们能够看到，当我们走向长距离或低能量时理论如何演化。它非常有用，因为它告诉我们哪些动力学性质在长距离或低能量下演生。
\item 由于一个有效理论只会演化成相似的有效理论，我们无法用重整化群方法得到性质上全新的现象的演生，例如从玻色模型演生出光和费米子。
\end{keypoints}

本节中，为了理解非紧致时钟模型的物理性质，我们将直接处理时钟模型中的 $\theta$ 场。

我们注意到，如果不同位置处的涨落 $O(\pmb{x})$ 与 $O(\pmb{y})$ 彼此独立地涨落，那么所谓的\emph{连通关联}（connected correlation）
\begin{equation*}
G_{\mathrm{conn}} = \langle O(\pmb{x})O(\pmb{y})\rangle - \langle O(\pmb{x})\rangle\langle O(\pmb{y})\rangle
\end{equation*}

% ---- 书页 111 (p126) ----
为零。所以，连通关联度量的正是 $O(\pmb{x})$ 与 $O(\pmb{y})$ 的涨落之间的关联。

当 $g_{l} = 0$ 时，无论 $\kappa_{l}$ 取何值，非紧致时钟模型总是具有代数长程关联：$\langle \mathrm{e}^{\mathrm{i}n\theta(\pmb{x})}\mathrm{e}^{-\mathrm{i}n\theta(0)}\rangle - \langle \mathrm{e}^{\mathrm{i}n\theta(\pmb{x})}\rangle\langle \mathrm{e}^{-\mathrm{i}n\theta(0)}\rangle \sim \frac{1}{|\pmb{x}|^{n^{2}/2\pi\kappa_{l}}}$。这里的问题是，$g_{l}\cos(n\theta)$ 项会如何影响代数长程关联。

如上一节所讨论，当 $\kappa_{l} < n^{2}/8\pi$ 时，$\mathrm{e}^{\mathrm{i}n\theta(\pmb{x})}$ 是不相关的，小的 $g_{l}\cos(n\theta)$ 项不会影响代数长程关联。当 $\kappa_{l} > n^{2}/8\pi$ 时，$\mathrm{e}^{\mathrm{i}n\theta(\pmb{x})}$ 是相关的。我们预期，无论 $\bar{g}_{l}$ 多么小，$g_{l}\cos(n\theta)$ 项都会把代数长程关联变为短程关联。我们看到，对于小的 $\bar{g}_{l}$，非紧致时钟模型在 $\kappa_{l} = n^{2}/8\pi$ 处有一个相变。下面我们将用 RG 方法来理解这个相变。

我们注意到，只有在指定了短距离截断 $l$ 之后，时钟模型才是良定义的。RG 方法的关键一步，是把波长介于 $l$ 与 $\lambda$（$\lambda > l$）之间的涨落积掉。这得到一个具有新截断 $\lambda$ 的模型。为了把短波长 $\theta$ 涨落积掉，我们先写
\begin{equation*}
\theta_{l} = \theta_{\lambda} + \delta\theta
\end{equation*}
其中 $\delta\theta$ 只包含波长介于 $l$ 与 $\lambda$ 之间的涨落。由于短波长涨落 $\delta\theta$ 受到 $\kappa(\partial_{\pmb{x}}\delta\theta)^{2}$ 项的压制，我们预期 $\delta\theta$ 很小，并把作用量展开到 $\delta\theta$ 的二阶：
\begin{equation*}
\begin{aligned}
S = {} & \int \mathrm{d}^{2}\pmb{x}\ \left( \frac{\kappa_{l}}{2}(\partial_{\pmb{x}}\theta_{\lambda})^{2} - g_{l}\cos(n\theta_{\lambda}) + \frac{\kappa_{l}}{2}(\partial_{\pmb{x}}\delta\theta)^{2} \right) \\
& + \int \mathrm{d}^{2}\pmb{x}\ \left( -ng_{l}\sin(\theta_{\lambda})\delta\theta + \frac{n^{2}}{2}g_{l}\cos(n\theta_{\lambda})(\delta\theta)^{2} \right)
\end{aligned}
\end{equation*}
（译注：原文 $\sin(\theta_{\lambda})$ 按展开式应为 $\sin(n\theta_{\lambda})$。）

我们把 $\theta_{\lambda}$ 当作光滑的背景场，并把 $\delta\theta$ 积掉（这种方法称为背景场 RG 方法）。我们得到有效作用量
\begin{equation*}
\begin{aligned}
S = {} & \int \mathrm{d}^{2}\pmb{x}\ \left( \frac{\kappa_{l}}{2}(\partial_{\pmb{x}}\theta_{\lambda})^{2} - g_{l}\cos(n\theta_{\lambda}) + \frac{1}{2}g_{l}\cos(n\theta_{\lambda})K(0) \right) \\
& - \int \mathrm{d}^{2}\pmb{x}\,\mathrm{d}^{2}\pmb{y}\ \frac{1}{2}(g_{l})^{2}\sin(n\theta_{\lambda}(\pmb{x}))K(\pmb{x}-\pmb{y})\sin(n\theta_{\lambda}(\pmb{y}))
\end{aligned}
\end{equation*}
